% ===================================================================
%  QUADRO N.º 3 — Infância e Juventude.
%
%  Ceci n'est PAS une transcription : c'est une traduction, faite en
%  2026, en portugais d'aujourd'hui. Voir texto/pt/10-quadro-01.tex
%  pour la consigne, qui vaut pour les seize fichiers.
%
%  « julho » n'est pas en gras ici, non plus qu'en ido : le fac-simile
%  l'oublie entre « junho » et « agosto », et c'est gravuri/korekti.json
%  qui le retablit, colonne par colonne, sur la page de lecture.
% ===================================================================

%%K t03-apar-1 apar
\VUsaut{3.0mm}
\VUtitre{15.0pt}{60}{QUADRO N.º 3}
\VUsaut{1.6mm}
\VUtitre{11.4pt}{0}{Infância e Juventude.}
\VUsaut{3.4mm}

%%K t03-apar-2 apar
(Os quadros 3 e 4 estão dispostos de modo que apresentam, em quatro
\VUgras{cenas familiares}, o vocabulário essencial do \VUgras{tempo}, da
\VUgras{idade}, da \VUgras{família}, da \VUgras{roupa} e da \VUgras{saúde}.)
\VUsaut{4.0mm}

%%K t03-tit-1 sub
\VUcentre{12.0pt}{0}{\textit{Primeira cena.}}
\VUsaut{3.0mm}

%%K t03-tit-2 sub
\VUpk{10.2pt}{40}{Um batizado na aldeia.}
\VUsaut{2.4mm}

%%K t03-tit-3 sub
\VUpk{10.2pt}{40}{A primavera.}
\VUsaut{3.0mm}

%%K t03-c1-01-1 p
1. --- Estamos na \VUgras{primavera}, no mês de \VUgras{março},
\VUgras{abril} ou \VUgras{maio}, num dia da \VUgras{semana} ---
\VUgras{segunda-feira}, \VUgras{terça-feira} ou \VUgras{quarta-feira}. São
\VUgras{dez} horas da manhã.

%%K t03-c1-02-1 p
2. --- Está bom \VUgras{tempo}, o \VUgras{ar} é puro. O sol brilha. Umas
\VUgras{nuvenzinhas} \textsuperscript{(2)} sobem pelo \VUgras{céu}
\textsuperscript{(1)} azul.

%%K t03-c1-03-1 p
3. --- As \VUgras{árvores} mal têm \VUgras{folhas}. Os esbeltos
\VUgras{álamos} \textsuperscript{(3)} e o alto \VUgras{plátano} \textsuperscript{(17)}
não trazem ainda mais do que rebentos.

%%K t03-c1-04-1 p
4. --- As \VUgras{andorinhas} \textsuperscript{(4)} voltaram e esvoaçam com
alegres chilreios à volta da \VUgras{flecha} \textsuperscript{(7)} do
\VUgras{campanário} \textsuperscript{(5)} que coroa a \VUgras{igreja}
\textsuperscript{(6)} da aldeia.

%%K t03-tit-4 sub
\VUsaut{3.4mm}
\VUpk{10.2pt}{40}{A Igreja.}
\VUsaut{3.0mm}

%%K t03-c1-05-1 p
5. --- Muito simples é esta igreja com o seu grande \VUgras{portal}
\textsuperscript{(13)} e o seu grosso \VUgras{relógio} \textsuperscript{(8)}. O
\VUgras{ponteiro pequeno} \textsuperscript{(9)} está sobre o algarismo X: marca
as \VUgras{horas}. O \VUgras{grande} \textsuperscript{(10)} está sobre o XII
(meio-dia); marca os \VUgras{minutos}.

%%K t03-c1-06-1 p
6. --- No campanário, os \VUgras{sinos} \textsuperscript{(11)} soam alegremente.
No alto do telhado ergue-se uma pequena \VUgras{cruz} \textsuperscript{(12)} de
pedra.

%%K t03-c1-07-1 p
7. --- A igreja não tem só um grande \VUgras{portal} \textsuperscript{(13)}, tem
também uma portinha lateral. Por esta porta o senhor \VUgras{pároco}
\textsuperscript{(15)}, vestido com a sua \VUgras{batina}, sai da
\VUgras{sacristia} \textsuperscript{(14)} e dirige-se à \VUgras{casa paroquial}
\textsuperscript{(19)}.

%%K t03-c1-08-1 p
8. --- Ao lado dele está a \VUgras{leiteira} \textsuperscript{(24)} com o seu
\VUgras{burro} \textsuperscript{(23)}. Volta da cidade, aonde levou o seu bom
leite para as crianças.

%%K t03-tit-5 sub
\VUsaut{4.0mm}
\VUpk{10.2pt}{40}{O pomar.}
\VUsaut{3.4mm}

%%K t03-c1-09-1 p
9. --- O senhor Aliborão (o burro) está muito contente com esta corrida
matinal. Aspira com deleite o ar embalsamado pelo perfume das
\VUgras{flores} \textsuperscript{(20)} que cobrem as \VUgras{árvores de fruto}
\textsuperscript{(18)} do \VUgras{pomar} vizinho. Todos rosados e brancos estão
esses \VUgras{pessegueiros} \textsuperscript{(18)} e esses \VUgras{damasqueiros}
\textsuperscript{(19)}, que prometem uma boa colheita.

%%K t03-c1-10-1 p
10. --- Entretanto, as pequenas \VUgras{abelhas} \textsuperscript{(22)} da
\VUgras{colmeia} \textsuperscript{(21)} vão lá libar o seu doce \VUgras{mel},
zumbindo.

%%K t03-c1-11-1 p
11. --- Mas o que se passa? As crianças da aldeia chegam a
correr e fazem fugir os \VUgras{gansos} \textsuperscript{(25)} assustados. É o
cortejo que sai da igreja depois do \VUgras{batizado} do filho do
notário.

%%K t03-c1-12-1 p
12. --- O pequeno Jaime, no seu \VUgras{berço} \textsuperscript{(26)}, acorda com
o alegre som dos sinos, e a sua irmã Madalena, que também quer ver o
cortejo, pede à mãe que venha penteá-la e vesti-la à soleira da casa.

%%K t03-c1-13-1 p
13. --- A mãe acede. Põe o seu \VUgras{lenço de cabeça} \textsuperscript{(28)}, o
seu \VUgras{corpete} \textsuperscript{(29)} e o seu \VUgras{avental}
\textsuperscript{(30)}, e vem sentar-se numa cadeirinha diante da porta.
Madalena traz apenas as suas \VUgras{anáguas} \textsuperscript{(31)} e o seu
\VUgras{espartilho} \textsuperscript{(32)}.

%%K t03-c1-14-1 p
14. --- Como é feliz! Esquece-se das suas flores preferidas, os seus
\VUgras{cravos} \textsuperscript{(34)}, onde pousam as \VUgras{borboletas}
\textsuperscript{(35)}, as suas \VUgras{tulipas} \textsuperscript{(36)}, os seus
\VUgras{lírios-do-vale} \textsuperscript{(37)} em \VUgras{vasos}
\textsuperscript{(38)} sobre o \VUgras{muro} \textsuperscript{(33)}, e as suas
\VUgras{rosas} \textsuperscript{(39)}, que formam uma sebe tão bela. Contempla
os ricos vestidos das senhoras do cortejo.

%%K t03-c1-15-1 p
15. --- Os rapazes da aldeia não fazem muito caso dos vestidos.
Precipitam-se e empurram-se para conseguir \VUgras{rebuçados}
\textsuperscript{(55)} ou \VUgras{moedas} \textsuperscript{(52)}. Um deles chora
\textsuperscript{(40)}.

%%K t03-c1-16-1 p
16. --- Outro pisou a pata do \VUgras{cão} \textsuperscript{(42)}, que foge a
ganir e faz fugir a \VUgras{galinha} \textsuperscript{(43)} e os seus
\VUgras{pintainhos} \textsuperscript{(44)}.

%%K t03-tit-6 sub
\VUsaut{5.0mm}
\VUpk{10.2pt}{40}{O cortejo do batizado.}
\VUsaut{4.0mm}

%%K t03-c1-17-1 p
17. --- Diante do cortejo caminha a \VUgras{ama} \textsuperscript{(45)}. Traz uma
bela \VUgras{touca} \textsuperscript{(46)} com fitas compridas. Leva ao colo o
\VUgras{recém-nascido} \textsuperscript{(47)}, todo vestido de \VUgras{rendas}
\textsuperscript{(48)}.

%%K t03-c1-18-1 p
18. --- Atrás da ama vêm o \VUgras{padrinho} \textsuperscript{(49)} e a
\VUgras{madrinha} \textsuperscript{(53)}. Distribuem os rebuçados e as moedas. O
padrinho traz \VUgras{luvas} \textsuperscript{(51)} e \VUgras{cartola}
\textsuperscript{(50)}. A madrinha traz na mão um belo \VUgras{ramo}
\textsuperscript{(54)}.

%%K t03-c1-19-1 p
19. --- Atrás deles vemos o \VUgras{tio} \textsuperscript{(56)} e a \VUgras{tia}
\textsuperscript{(58)} da criança. A tia traz um elegante \VUgras{chapéu}
\textsuperscript{(59)}, o tio uma \VUgras{sobrecasaca} \textsuperscript{(57)} preta e
um colete branco. Vêm depois o \VUgras{primo} \textsuperscript{(60)} e a
\VUgras{prima} \textsuperscript{(61)}. Esta leva pela mão a \VUgras{irmã}
\textsuperscript{(62)} do recém-nascido, a pequena Berta.

%%K t03-c1-20-1 p
20. --- O outro \VUgras{irmão} \textsuperscript{(63)} de Berta caminha atrás. Dá a
mão a uma \VUgras{parenta} \textsuperscript{(64)}. Os \VUgras{amigos}
\textsuperscript{(65)} da família vêm a seguir.

%%K t03-tit-7 sub
\VUsaut{4.6mm}
\VUpk{10.2pt}{40}{Os que olham.}
\VUsaut{3.4mm}

%%K t03-c1-21-1 p
21. --- Perto do cortejo, um \VUgras{mendigo} \textsuperscript{(66)} apoiado
na sua \VUgras{muleta} \textsuperscript{(67)}. Pede esmola.

%%K t03-c1-22-1 p
22. --- Do outro lado há um menino muito \VUgras{cortês} \textsuperscript{(68)}.
Cumprimenta e dá os \VUgras{«bons dias»} às pessoas que conhece.

%%K t03-c1-23-1 p
23. --- Os da aldeia saíram de suas casas para ver passar o cortejo do
batizado. Vemos um \VUgras{camponês} \textsuperscript{(69)} com o seu
\VUgras{barrete de algodão} e ao seu lado uma \VUgras{camponesa}
\textsuperscript{(70)}. Depois uma \VUgras{velha} \textsuperscript{(71)} apoiada na sua
\VUgras{bengala} \textsuperscript{(72)}. Por fim, o \VUgras{moleiro}
\textsuperscript{(73)} com o seu largo \VUgras{chapéu de feltro}
\textsuperscript{(74)}, e uma camponesa jovem calçada com \VUgras{tamancos}
\textsuperscript{(75)}, que ensina a andar ao seu menino.

%%K t03-c1-24-1 p
24. --- É uma cena muito divertida.
\VUsaut{4.0mm}

%%K t03-tit-8 sub
\VUcentre{12.0pt}{0}{\textit{Segunda cena.}}
\VUsaut{3.0mm}

%%K t03-tit-9 sub
\VUpk{10.2pt}{40}{Uma festa pública.}
\VUsaut{2.4mm}

%%K t03-tit-10 sub
\VUpk{10.2pt}{40}{Jogos náuticos e regatas.}
\VUsaut{3.0mm}

%%K t03-c2-01-1 p
1. --- Chegou o \VUgras{verão}. O sol ardente do mês de \VUgras{junho}
(julho ou \VUgras{agosto}) convida os jovens a banhar-se e a sair de
barco.

%%K t03-c2-02-1 p
2. --- Na margem direita do rio, os remadores despem-se para tomar parte
nos jogos náuticos e nas regatas.

%%K t03-c2-03-1 p
3. --- Um deles tira os \VUgras{sapatos} \textsuperscript{(1)}; desata-os
(desabotoa-os). Os seus dois companheiros já estão prontos. Só trazem a
sua \VUgras{camisola} \textsuperscript{(2)} e o seu \VUgras{calção de banho}
\textsuperscript{(3)}. Conversam enquanto esperam os amigos. Falam da feira e
da festa que se celebra na outra margem.

%%K t03-c2-04-1 p
4. --- No chão, junto deles, jaz a \VUgras{roupa} dos seus companheiros:
um \VUgras{par} de \VUgras{botins} \textsuperscript{(4)}, \VUgras{sapatos}
\textsuperscript{(5)}, \VUgras{meias} \textsuperscript{(6)}, chapéus de
\VUgras{feltro} \textsuperscript{(7)} e de \VUgras{palha} \textsuperscript{(8)} e uma
\VUgras{gravata} \textsuperscript{(9)}.

%%K t03-c2-05-1 p
5. --- Um dos jovens dobrou com cuidado o seu \VUgras{casaco}
\textsuperscript{(10)}. Põe-no sobre uma toalha, na qual o seu vizinho
embrulhará já a seguir os dois \VUgras{fatos}.

%%K t03-c2-06-1 p
6. --- Um sexto rapaz vai atrasado. Traz ainda o seu \VUgras{boné}
\textsuperscript{(11)} na cabeça. Não tirou nem a \VUgras{camisa}
\textsuperscript{(12)}, nem o \VUgras{colarinho} \textsuperscript{(13)}, nem os
\VUgras{punhos} \textsuperscript{(14)}. Tem na mão o seu \VUgras{colete}
\textsuperscript{(17)} e traz ainda as suas \VUgras{calças} \textsuperscript{(16)} e os
seus \VUgras{suspensórios} \textsuperscript{(15)}. De certeza que não estará
pronto para a próxima corrida.

%%K t03-c2-07-1 p
7. --- Perto dos jovens, um carreiro orlado de \VUgras{cardos}
\textsuperscript{(18)} desce à margem do rio, onde o tio Jaime prepara, entre
os \VUgras{juncos} \textsuperscript{(19)}, o \VUgras{fogo de artifício}
\textsuperscript{(20)} que se lançará à noite.

%%K t03-tit-11 sub
\VUsaut{4.4mm}
\VUpk{10.2pt}{40}{A corrida.}
\VUsaut{3.4mm}

%%K t03-c2-08-1 p
8. --- No rio, algumas senhoras que se resguardam sob as suas
sombrinhas passeiam de \VUgras{barco} \textsuperscript{(21)}. Seguem a corrida,
que é muito interessante. Em cada \VUgras{iole} \textsuperscript{(22)}, quatro
\VUgras{remadores} \textsuperscript{(24)} remam com todas as suas forças,
enquanto o \VUgras{timoneiro} \textsuperscript{(26)} segura o \VUgras{leme}
\textsuperscript{(25)} e dirige a frágil embarcação. Os \VUgras{remos}
\textsuperscript{(23)} batem a água a compasso, e os ligeiros botes navegam a
uma velocidade vertiginosa.

%%K t03-c2-09-1 p
9. --- Alguns jovens disputam, junto ao pilar da ponte, o prémio do
\VUgras{pau-de-sebo horizontal} \textsuperscript{(28)}. Um deles avança com
cautela pelo pau flexível e escorregadio para alcançar a pequena
\VUgras{bandeira} \textsuperscript{(29)} presa na ponta. O seu desditoso
\VUgras{competidor} \textsuperscript{(30)} caiu à água. Nada para chegar à
margem.

%%K t03-c2-10-1 p
10. --- Uns rapazes mais velhos \VUgras{justam em barco} \textsuperscript{(32)}
junto ao \VUgras{cais} \textsuperscript{(33)}. A todos estes jogos assiste um
\VUgras{público} \textsuperscript{(31)} numeroso, apoiado no \VUgras{parapeito} da
\VUgras{ponte} \textsuperscript{(27)} e apinhado na \VUgras{margem}. Alguns
espectadores, porém, voltam-se para seguir o jogo do \VUgras{pau-de-sebo}
\textsuperscript{(45)} ou para admirar o \VUgras{cortejo municipal} que acode à
\VUgras{distribuição} dos prémios.

%%K t03-c2-11-1 p
11. --- Primeiro vêm uns \VUgras{garotos} \textsuperscript{(35)} que vão diante
dos \VUgras{músicos} \textsuperscript{(36)}; depois vêm o \VUgras{presidente da
câmara} \textsuperscript{(37)}, os vereadores e a \VUgras{câmara municipal} da
pequena cidade. Fecham a marcha os \VUgras{bombeiros} \textsuperscript{(38)}.

%%K t03-tit-12 sub
\VUsaut{5.0mm}
\VUpk{10.2pt}{40}{A feira.}
\VUsaut{3.6mm}

%%K t03-c2-12-1 p
12. --- Há muita gente no \VUgras{recinto da feira} \textsuperscript{(39)}.
Vem-se ver as diferentes \VUgras{barracas}: os \VUgras{cavalinhos de
madeira} \textsuperscript{(40)}, o \VUgras{circo} \textsuperscript{(41)}, onde o
espetáculo já começou; o \VUgras{baile público} \textsuperscript{(42)}, onde os
jovens e as jovens da cidade gozam o prazer de \VUgras{dançar}.

%%K t03-c2-13-1 p
13. --- Ao fundo vemos a \VUgras{praça de touros} \textsuperscript{(44)}, onde o
valente \VUgras{matador} espanhol luta contra o \VUgras{touro}
\textsuperscript{(45)} furioso; depois a \VUgras{montanha-russa}
\textsuperscript{(49)} e a \VUgras{barraca de tiro} \textsuperscript{(46)}, onde se
exercitam em manejar as \VUgras{armas de fogo} e em atirar ao
\VUgras{alvo}.

%%K t03-c2-14-1 p
14. --- Perto dali fica o \VUgras{teatro de feira} \textsuperscript{(47)}, no qual
se representam a \VUgras{comédia} e o \VUgras{drama}. Muito perto está a
barraca do \VUgras{gigante} \textsuperscript{(48)} e do \VUgras{anão}. A
\VUgras{estatura} do primeiro é de dois metros e quinze; o segundo mal é
tão alto como uma criança.

%%K t03-c2-15-1 p
15. --- Por fim vem a grande \VUgras{coleção de feras} \textsuperscript{(50)} com
os seus \VUgras{animais selvagens}, o seu \VUgras{tigre} \textsuperscript{(51)} de
poderosas \VUgras{garras}, o seu \VUgras{leão} \textsuperscript{(52)} de espessa
\VUgras{juba}, as suas duas \VUgras{girafas} \textsuperscript{(53)} e o seu
\VUgras{elefante} \textsuperscript{(54)}, que usa com tanta manha a sua
\VUgras{tromba}.

%%K t03-c2-16-1 p
16. --- O \VUgras{domador} \textsuperscript{(55)} que amestra estes animais está à
porta da barraca.
\VUsaut{6.0mm}
\VUfilet{22mm}
